% answers_short.tex —— 2006 年·普通高等学校招生考试·重庆卷理科 · 简短答案速查
% 由 tools/make_short_answers.py 自 answers.tex 的 \gksolhead/\gkanswer 提取，
% 并与 2006/2006_数学真题_重庆理_答案要点.md 逐题比对后生成，勿手改（改源 MD 后重跑）。

\begin{gknotice}
\textbf{使用说明}：本节只给最终答案便于快速核对，完整推导见解析版。答案为编写者独立推导，\textbf{非官方标准答案}。
\end{gknotice}

\gkpart{选择题}
\begin{center}\begin{tabular}{ccccccccccc}
\toprule
题号 & 1 & 2 & 3 & 4 & 5 & 6 & 7 & 8 & 9 & 10 \\
\midrule
答案 & D & B & A & C & A & C & B & B & D & D \\
\bottomrule
\end{tabular}\end{center}

\gkpart{填空题}
\begin{itemize}[leftmargin=2.6em]
  \item[\textbf{11.}] (1+7i)/10。
  \item[\textbf{12.}] 1/2。
  \item[\textbf{13.}] $-$56/65。
  \item[\textbf{14.}] $a_n=2^{n+1}-3$
  \item[\textbf{15.}] (2,3)。
  \item[\textbf{16.}] a>1。
\end{itemize}

\gkpart{解答题}
\begin{itemize}[leftmargin=2.6em]
  \item[\textbf{17.}] (1) $\omega=\frac12$；(2) a=(√3+1)/2
  \item[\textbf{18.}] (1) $\xi\sim\operatorname{Bin}(5,\frac13)$，$P(\xi=k)=\binom{5}{k}(\frac13)^k(\frac23)^{5-k}$，$k=0,\ldots,5$；(2) $E\xi=\frac53$
  \item[\textbf{19.}] (1) CD$\perp$平面 BEF；(2) k>2/√15
  \item[\textbf{20.}] (1) U/PENDING；(2) $-$6≤b≤2
  \item[\textbf{21.}] (1) f(1)=1，f(a)=a；(2) f(x)=x²$-$x+1
  \item[\textbf{22.}] (1) $b_n\le\frac{\sqrt3}{2}$；(2) $S_1<S_2$ 且 $S_n>S_{n+1}$（$n\ge3$）
\end{itemize}
